\documentclass[11pt,a4paper]{article}
\usepackage[T1]{fontenc}
\usepackage{lmodern}
\usepackage{isabelle,isabellesym}
\usepackage{amsmath,amssymb,amsthm}
\usepackage{stmaryrd}
\usepackage{wasysym}
\usepackage{textcomp}
\usepackage[margin=1in]{geometry}
\usepackage{pdfsetup}

\isabellestyle{it}
\isadroptag{theory}
\renewcommand{\isastyletext}{\normalsize\normalfont}
\renewcommand{\isastyletxt}{\normalfont}

\newcommand{\ket}[1]{|#1\rangle}
\newcommand{\bra}[1]{\langle #1|}
\newcommand{\norm}[1]{\|#1\|}
\DeclareMathOperator{\Tr}{Tr}
\newcommand{\C}{\mathbb{C}}
\newcommand{\R}{\mathbb{R}}
\newcommand{\Hil}{\mathcal{H}}

\begin{document}

\title{Essential Functional Analysis\\[0.5em]
  \large Prerequisites for the formalisation of Vern Paulsen's lectures on quantum computing}
\author{Dustin Bryant}
\maketitle

\begin{abstract}
General facts about complex matrices that the formalised lectures (session
\texttt{Quantum\_Computing\_CP\_Maps}) build on: sums, adjoints, unitary matrices and Kronecker
products, linear-space isomorphisms, and matrices of fixed dimension as a type, with the
Hilbert--Schmidt norm (a normed space for the Hahn--Banach theorem of \texttt{HOL-Hahn\_Banach})
and the structure of a Euclidean space (so that \texttt{HOL-Analysis} applies).  Nothing here depends on the lectures.
\end{abstract}

\tableofcontents
\newpage

\parindent 0pt\parskip 0.5ex

\input{Matrix_Lemmas}
\input{Unitary_Tensor_Algebra}
\input{Matrix_Typed}
\input{Lmat_Euclidean}
\input{Essential_Functional_Analysis}

\end{document}
